\unitopen{01}
\begin{candos}
\item przeczytać na głos dowolne niemieckie słowo bez zgadywania (alfabet i reguły wymowy);
\item przywitać się i pożegnać o każdej porze dnia, formalnie (\de{Sie}) i nieformalnie (\de{du});
\item podziękować, przeprosić i zapytać \de{Wie geht's?} oraz odpowiedzieć;
\item literować swoje imię i nazwisko po niemiecku.
\end{candos}

\section{Słownictwo}
\vocabtable{01}

\section{Alfabet i nazwy liter}
Niemiecki alfabet ma 26 liter jak łaciński, plus \de{Ä}, \de{Ö}, \de{Ü} i \de{ß}. Nazwy liter
potrzebne Ci do literowania nazwiska przez telefon (tak robi się w Szwajcarii i Niemczech przy
każdej rezerwacji i rejestracji):

\begin{center}\small
\begin{tabular}{@{}*{9}{l}@{}}
A [a] & B [be] & C [ce] & D [de] & E [e] & F [ef] & G [ge] & H [ha] & I [i]\\
J [jot] & K [ka] & L [el] & M [em] & N [en] & O [o] & P [pe] & Q [ku] & R [er]\\
S [es] & T [te] & U [u] & V [fał] & W [we] & X [iks] & Y [ypsilon] & Z [cet] & \\
\end{tabular}
\end{center}

\de{Ä} (\de{A-Umlaut}), \de{Ö}, \de{Ü} i \de{ß} (\de{Eszett} albo \de{scharfes S}) to osobne dźwięki,
nie ozdobniki. Konrad Cinkusz literujesz: \de{Ka-O-En-Er-A-De}. W polskim imieniu \de{ł} i \de{ń} nie
istnieją: przy literowaniu podaj \de{L} oraz \de{N} i powiedz, że to polskie litery.

\section{Wymowa: reguły, które pokrywają 90\% słów}
Niemiecki jest dużo bardziej regularny niż angielski: zapis prawie zawsze mówi, jak czytać.

\begin{center}\small
\begin{tabular}{@{}lll@{}}
\toprule
\textsf{\textbf{Zapis}} & \textsf{\textbf{Czytasz jak (przybliżenie)}} & \textsf{\textbf{Przykład}}\\
\midrule
\de{ei} & \emph{aj} & \de{mein} [majn], \de{Zeit} [cajt]\\
\de{ie} & długie \emph{i} & \de{sieben} [zibn], \de{Bier}\\
\de{eu}, \de{äu} & \emph{oj} & \de{Deutsch} [dojcz], \de{Häuser}\\
\de{sch} & \emph{sz} & \de{Schule} [szule], \de{schnell}\\
\de{sp-}, \de{st-} na początku słowa & \emph{szp-}, \emph{szt-} & \de{Sprache} [szprache], \de{Stadt}\\
\de{z}, \de{tz} & \emph{c} & \de{zehn} [cen], \de{Platz}\\
\de{w} & jak polskie \emph{w} (nie angielskie) & \de{Wasser}, \de{wohnen}\\
\de{v} & zwykle \emph{f} & \de{Vater} [fater], \de{viel} [fil]\\
\de{s} przed samogłoską & \emph{z} (dźwięczne) & \de{Sonne} [zone], \de{sieben}\\
\de{ß}, \de{ss} & \emph{s} (długie / krótkie) & \de{Straße}, \de{muss}\\
\de{j} & jak polskie \emph{j} & \de{ja}, \de{Jahr}\\
\de{ä} & jak \emph{e} (długie lub krótkie) & \de{Mädchen}, \de{Käse}\\
\de{ö} & \emph{e} z zaokrąglonymi wargami & \de{schön}, \de{Köln}\\
\de{ü} & \emph{i} z zaokrąglonymi wargami & \de{Tür}, \de{über}\\
\bottomrule
\end{tabular}
\end{center}

\subsection*{Dwa \de{ch}}
Po \de{a}, \de{o}, \de{u}, \de{au} czytasz twarde \de{ch} jak polskie \emph{ch}: \de{Buch}, \de{Nacht}, \de{Bach}.
Po pozostałych samogłoskach i po \de{l}, \de{n}, \de{r} czytasz miękkie \de{ch} (jak szept polskiego \emph{ś},
bez wibracji): \de{ich}, \de{Mädchen}, \de{Milch}, \de{Mönch}.

\subsection*{Długość samogłosek}
Przed podwójną spółgłoską samogłoska jest krótka (\de{Bett}, \de{Mutter}); po samogłosce \de{h} jest ona
długa i \de{h} nie brzmi (\de{gehen}, \de{Zahl}).

\subsection*{Ostatnia spółgłoska traci dźwięczność}
Na końcu słowa \de{b}, \de{d}, \de{g} czytasz jak \emph{p}, \emph{t}, \emph{k}: \de{Tag} [tak], \de{und} [unt],
\de{Hund} [hunt]. Polacy mają to w języku i nie mają problemu.

\begin{falle}
Polskie \emph{rz} i \emph{ż} nie istnieją po niemiecku. \de{Jazz} to [dżez] tylko w zapożyczeniu.
Twoja największa pułapka to \textbf{\de{ch} w \de{ich}}: czytane jak polskie \emph{ch} brzmi jak \emph{ik}
z chrypką. Druga: \de{w} i \de{v} --- nie zamieniaj ich na angielskie.
\end{falle}

\section{Powitania: formalnie i nieformalnie}
Niemiecki rozróżnia \de{Sie} (Pan/Pani/Państwo; zawsze wielką literą) i \de{du} (ty). W biurze i na
początku znajomości z nieznajomym używasz \de{Sie}.

\begin{center}\small
\begin{tabular}{@{}lll@{}}
\toprule
 & \textsf{\textbf{Formalnie (Sie)}} & \textsf{\textbf{Nieformalnie (du)}}\\
\midrule
Powitanie & \de{Guten Tag}, \de{Guten Morgen}, \de{Guten Abend} & \de{Hallo}, \de{Guten Morgen}\\
Pożegnanie & \de{Auf Wiedersehen} & \de{Tschüss}, \de{Bis bald}\\
Pytanie o samopoczucie & \de{Wie geht es Ihnen?} & \de{Wie geht's?}, \de{Wie geht es dir?}\\
Odpowiedź & \de{Danke, gut. Und Ihnen?} & \de{Gut, danke. Und dir?}\\
\bottomrule
\end{tabular}
\end{center}

\begin{szwajcaria}
W szwajcarskim biurze zaczniesz od \de{Grüezi} (dialektalne, ale zrozumiałe wszędzie) albo
\de{Guten Tag}. W IT i start-upach w Zurychu kolegów zwykle się tykać (\de{du}), nawet szefa; w kontakcie
z klientem, urzędem i podczas rozmowy kwalifikacyjnej używaj \de{Sie}, dopóki ktoś nie zaproponuje inaczej.
Zaproponować \de{du} (\de{Wir können uns duzen}) wypada starszej osobie albo wyższej rangą, nie Tobie.
Szwajcarzy mówią \de{Merci} zamiast \de{Danke} i \de{Ade} / \de{Uf Wiederluege} przy pożegnaniu.
\end{szwajcaria}

\begin{praca}
Na pierwszym spotkaniu online powiedz: \de{Guten Morgen zusammen!} (do wszystkich), a po spotkaniu
\de{Danke und einen schönen Tag noch!} Pamiętaj: \de{Guten Morgen} do około 10--11, \de{Guten Tag} potem,
\de{Guten Abend} od późnego popołudnia.
\end{praca}

\section{Dialogi}
\textbf{Dialog 1 --- w recepcji (formalnie)}
\begin{dialog}
\turn{A}{Guten Morgen! Wie geht es Ihnen?}{Dzień dobry! Jak się Pan/Pani miewa?}
\turn{B}{Guten Morgen! Danke, gut. Und Ihnen?}{Dzień dobry! Dziękuję, dobrze. A Panu/Pani?}
\turn{A}{Auch gut, danke. Willkommen!}{Też dobrze, dziękuję. Witamy!}
\turn{B}{Vielen Dank. Auf Wiedersehen!}{Bardzo dziękuję. Do widzenia!}
\end{dialog}

\textbf{Dialog 2 --- koledzy (nieformalnie)}
\begin{dialog}
\turn{A}{Hallo! Wie geht's?}{Cześć! Jak się masz?}
\turn{B}{Hallo! Gut, danke. Und dir?}{Cześć! Dobrze, dziękuję. A tobie?}
\turn{A}{Auch gut. Bis morgen!}{Też dobrze. Do jutra!}
\turn{B}{Tschüss!}{Cześć!}
\end{dialog}

\section{Ćwiczenia}
\begin{cwiczenia}
\exitem{Jak czytasz \de{Zeit}? Zapisz polskimi literami.}{cajt (\de{ei} = \emph{aj}, \de{z} = \emph{c})}
\exitem{Jak czytasz \de{Schule}?}{szule}
\exitem{Jak czytasz \de{Vater}?}{fater (\de{v} = \emph{f})}
\exitem{Jak czytasz \de{sieben}?}{zibn / ziben (\de{ie} = długie \emph{i}, \de{s} przed samogłoską = \emph{z})}
\exitem{Które \de{ch} jest twarde: \de{Buch} czy \de{ich}?}{\de{Buch} (po \de{u}); \de{ich} ma miękkie \de{ch}}
\exitem{Napisz po niemiecku: „Dzień dobry” o 15:00, formalnie.}{\de{Guten Tag}}
\exitem{Jak powiesz „Do widzenia” w rozmowie z szefem klienta?}{\de{Auf Wiedersehen}}
\exitem{Który zwrot jest nieformalny: \de{Wie geht's?} czy \de{Wie geht es Ihnen?}}{\de{Wie geht's?}}
\exitem{Odpowiedz formalnie na \de{Wie geht es Ihnen?}}{\de{Danke, gut. Und Ihnen?}}
\exitem{Przetłumacz: „Przepraszam, gdzie jest toaleta?” --- tylko pierwsze słowo.}{\de{Entschuldigung}}
\exitem{Przeliteruj \de{Hallo} po niemiecku (nazwy liter).}{\de{Ha-A-El-El-O} (ha, a, el, el, o)}
\exitem{Dlaczego \de{Tag} czytamy [tak], nie [tag]?}{Na końcu słowa \de{g} traci dźwięczność i brzmi jak \emph{k}}
\end{cwiczenia}

\section{Trening słownictwa (generowany)}
\drill{01}

\weekend{Nagraj (telefonem) 60 sekund: powitaj się formalnie, potem nieformalnie, podziękuj, przeproś,
przeliteruj swoje imię i nazwisko. Posłuchaj nagrania i porównaj z lektorem (DW \emph{Nicos Weg}, lekcja~1
albo dowolny nagrany dialog). Wyślij nagranie do AI lub korepetytora do korekty wymowy
\de{ch}, \de{sch}, \de{w}.}
